\chapter{Transmission Lines, Reflection, and SWR}
\label{ch:lines}
Every circuit so far has been \emph{lumped}: a finite set of components, a finite
state vector, and the assumption from \cref{ch:foundations} that the circuit is
small compared with a wavelength. A feed line breaks that assumption, and this
chapter is what replaces it. The reward is that almost nothing new is needed---the
line turns out to be built from the same \(L\) and \(C\) element laws, obeys the
same KVL and KCL, and produces the same resonances as \cref{ch:rlc,ch:rlcpar}, only
distributed along a length instead of concentrated in two parts.

The development is staged. The line is treated as \emph{lossless} from here through
\cref{sec:sparams}---the right first model, and the one that makes every result
clean. \Cref{sec:lineloss} then puts the loss back and shows what changes: feed-line
attenuation, why it climbs with frequency, and why a mismatch measured at the
transmitter looks better than it is. Everything after that---input impedance, stubs,
transformers, the Smith chart and the antenna---returns to the lossless model, because
loss changes the numbers without changing the structure.

\section{Why Lines Matter}
Transmission lines are distributed systems. When a feed line is a significant
fraction of a wavelength, voltage and current vary along the line and a lumped
wire approximation fails.

\begin{controlsbox}
A transmission line is a distributed-parameter dynamical system rather than a
lumped ODE. A lumped RLC circuit has a finite state vector; a line has a
continuum of states \(v(x,t)\) and \(i(x,t)\) indexed by position \(x\). Position
becomes part of the state description, so the model is a partial differential
equation with an infinite set of natural modes rather than a handful of
eigenvalues.

This is the edge \cref{sec:infinitepoles} walked up to: order \(n\) needs no new
machinery for any finite \(n\), but the \(n\to\infty\) limit of a ladder of
identical sections is qualitatively different---it is a delay, and a delay has
infinitely many crossings of every phase condition. The transition is continuous,
which is why so much of this chapter is Part~III's algebra reused.
\end{controlsbox}

\section{The Line as a Distributed System: Telegrapher's Equations}
Model an infinitesimal slice \(\dd x\) of a lossless line as a series inductance
\(L'\,\dd x\) and a shunt capacitance \(C'\,\dd x\), where \(L'\) and \(C'\) are
inductance and capacitance per unit length. Applying KVL and KCL to the slice
gives the \emph{telegrapher's equations},
\[
\frac{\partial v}{\partial x}=-L'\frac{\partial i}{\partial t},
\qquad
\frac{\partial i}{\partial x}=-C'\frac{\partial v}{\partial t}.
\]
These are the \emph{same} two element laws used throughout Part~III---the inductor's
\(v=L\,\dd i/\dd t\) and the capacitor's \(i=C\,\dd v/\dd t\)---applied to a slice
rather than to a discrete part. And notice their structure: a coupled pair of
first-order equations, one per energy store, exactly as in \cref{ch:rlc}. The only
difference is that the derivatives on the left are with respect to \emph{position}
instead of time, because the ``next'' store is the neighboring slice rather than a
separate component.

\subsection{From the Coupled Pair to the Wave Equation}
The series RLC chapter reduced its coupled pair to one second-order equation by
differentiating and substituting (\cref{ch:rlc}). Do the same here. Differentiate
the first telegrapher's equation with respect to \(x\):
\[
\frac{\partial^2 v}{\partial x^2}
=-L'\frac{\partial}{\partial x}\frac{\partial i}{\partial t}
=-L'\frac{\partial}{\partial t}\frac{\partial i}{\partial x},
\]
swapping the order of the two derivatives. Now substitute the second telegrapher's
equation, \(\partial i/\partial x=-C'\,\partial v/\partial t\):
\[
\frac{\partial^2 v}{\partial x^2}
=-L'\frac{\partial}{\partial t}\Bigl(-C'\frac{\partial v}{\partial t}\Bigr)
\qquad\Longrightarrow\qquad
\boxed{\;\frac{\partial^2 v}{\partial x^2}=L'C'\,\frac{\partial^2 v}{\partial t^2}.\;}
\]
This is the \emph{wave equation}. It is the distributed analogue of the boxed
second-order ODE of \cref{ch:rlc}, reached by the identical manoeuvre.

Any function of the combination \(t\mp x/v_p\) satisfies it, provided
\[
v_p=\frac{1}{\sqrt{L'C'}},
\]
which is quickest to check by substitution: for \(v=f(t-x/v_p)\), each
\(\partial/\partial x\) contributes a factor \(-1/v_p\) and each
\(\partial/\partial t\) a factor \(1\), so the equation reduces to
\(f''/v_p^2=L'C'f''\). Solutions are therefore \emph{traveling waves}, moving in
\(+x\) or \(-x\) at speed \(v_p\) without changing shape.

\subsection{Where \texorpdfstring{\(Z_0\)}{Z0} Comes From}
The characteristic impedance is not a separate assumption; it follows from the same
substitution. Take a single forward wave \(v=f(t-x/v_p)\) and put it into the first
telegrapher's equation:
\[
\underbrace{-\frac{1}{v_p}f'}_{\partial v/\partial x}
=-L'\frac{\partial i}{\partial t}
\qquad\Longrightarrow\qquad
\frac{\partial i}{\partial t}=\frac{f'}{L'v_p}
\qquad\Longrightarrow\qquad
i=\frac{f}{L'v_p},
\]
integrating in time. The ratio of voltage to current in that wave is therefore a
constant, independent of shape and of frequency:
\[
\frac{v}{i}=L'v_p=\frac{L'}{\sqrt{L'C'}}
=\boxed{\;Z_0=\sqrt{\frac{L'}{C'}}\;}
\]
So \(Z_0\) and the velocity factor are not empirical constants---they fall out of
the per-unit-length geometry of the line. Note also what \(Z_0\) is \emph{not}: it
dissipates nothing. It is a ratio maintained by a wave in flight, which is why a
resistor equal to \(Z_0\) can absorb a wave perfectly (it presents the same ratio
the wave already has) while the lossless line itself consumes no power.

\section{Wavelength and Velocity Factor}
The \textbf{velocity factor} \(VF=v_p/c\) is the ratio of the wave's speed in the
line to the speed of light in vacuum---a dimensionless number less than one. With it,
the wavelength \emph{on the line} at frequency \(f\) is
\[
\lambda=\frac{v_p}{f}=\frac{VF\,c}{f}.
\]
Because \(v_p=1/\sqrt{L'C'}\) and a dielectric raises \(C'\), a line
with dielectric has \(VF<1\); its physical length is shorter than the free-space
electrical length. That can be closed into a formula rather than left as a tendency.
A non-magnetic insulator leaves \(L'\) alone but multiplies \(C'\) by the relative
permittivity \(\varepsilon_r\), so
\[
v_p=\frac{1}{\sqrt{L'\varepsilon_rC'_0}}=\frac{c}{\sqrt{\varepsilon_r}}
\qquad\Longrightarrow\qquad
\boxed{\;VF=\frac{1}{\sqrt{\varepsilon_r}}\;}
\]
where \(C'_0\) is the capacitance the same geometry would have in vacuum. This is the
same \(v_p=1/\sqrt{L'C'}\) already derived, with the dielectric's effect named---and
it predicts the table below rather than merely being consistent with it.
Polyethylene has \(\varepsilon_r\approx2.25\), giving \(VF=1/1.5=0.667\); PTFE's
\(2.1\) gives \(0.69\); a foam that is half air gives about \(0.82\). Compare those
with the measured values quoted next. Typical values span the dielectric fill: solid-polyethylene
coax (RG-8, RG-58) runs \(VF\approx0.66\); foam-dielectric coax
\(VF\approx0.78\text{--}0.85\); PTFE coax \(VF\approx0.70\); flat twin-lead
\(VF\approx0.82\); and mostly-air open-wire or ``ladder'' line
\(VF\approx0.90\text{--}0.95\). Overhead power lines, essentially air-insulated,
are near \(VF\approx0.95\text{--}0.98\). Always use the actual line's \(VF\) when
cutting stubs or matching sections to length.

\section{Characteristic Impedance}
\(Z_0=\sqrt{L'/C'}\) is the voltage-to-current ratio of a traveling wave. Common
coaxial systems use \(\SI{50}{\ohm}\); \(\SI{75}{\ohm}\), \(\SI{300}{\ohm}\), and
\(\SI{450}{\ohm}\) appear in specific applications. A line terminated in \(Z_0\)
looks infinitely long---no reflection---because the load absorbs the traveling
wave exactly as more line would.

\section{Reflection Coefficient}
\label{sec:reflectionsection}
A wave traveling toward a load carries a fixed voltage-to-current ratio \(Z_0\). The
load demands its own ratio, \(Z_L\). If the two disagree, one wave cannot satisfy
both---so a second, backward-traveling wave appears, sized precisely to reconcile
them. That is all a reflection is, and the reconciliation is a two-line calculation.

Write the total voltage and current at the load as the sum of a forward wave
\(V^{+}\) and a backward wave \(V^{-}\). Voltages add; currents \emph{subtract},
because the backward wave carries its power the other way, so its current reverses
sign relative to its voltage:
\[
V=V^{+}+V^{-},
\qquad
I=\frac{V^{+}}{Z_0}-\frac{V^{-}}{Z_0}.
\]
The load imposes one boundary condition, \(V/I=Z_L\):
\[
Z_L=Z_0\,\frac{V^{+}+V^{-}}{V^{+}-V^{-}} .
\]
Solving for the ratio of reflected to incident amplitude defines the
\textbf{reflection coefficient} \(\Gamma=V^{-}/V^{+}\):
\[
Z_L\bigl(V^{+}-V^{-}\bigr)=Z_0\bigl(V^{+}+V^{-}\bigr)
\;\Rightarrow\;
V^{-}(Z_L+Z_0)=V^{+}(Z_L-Z_0)
\;\Rightarrow\;
\boxed{\;\Gamma=\frac{V^{-}}{V^{+}}=\frac{Z_L-Z_0}{Z_L+Z_0}.\;}
\]
Every limiting case now follows by inspection rather than by memory. If
\(Z_L=Z_0\) the numerator vanishes, \(\Gamma=0\): the load already presents the
ratio the wave carries, so nothing reflects. If the line is \emph{shorted}
(\(Z_L=0\)) then \(\Gamma=-1\)---full reflection with the voltage inverted, since the
short forces \(V=0\). If it is \emph{open} (\(Z_L\to\infty\)) then \(\Gamma=+1\):
full reflection in phase, since the open forces \(I=0\). Both extremes give
\(|\Gamma|=1\), as they must---a lossless termination that absorbs nothing has to
return everything.

\begin{controlsbox}
This is the same conjugate-of-a-mismatch bookkeeping as \cref{ch:resistive}'s
maximum-power-transfer condition, seen dynamically. There, a source delivered the
most power when the load matched it; here, a wave transfers all of its power when the
load matches \(Z_0\). ``Matched'' means the same thing in both places: the load is
presenting the ratio the source side is built to work into.
\end{controlsbox}

\section{Wave Coordinates: S-Parameters}
\label{sec:sparams}
The previous section did the hard part. Once \(V\) and \(I\) have been split into a
forward and a backward wave, the standard language of RF measurement is available for
the price of a change of variables---and the change of variables is the same kind of
move the state-space chapters keep making.

Define, at any port,
\[
a=\frac{V+Z_0I}{2\sqrt{Z_0}},
\qquad
b=\frac{V-Z_0I}{2\sqrt{Z_0}} .
\]
Substitute the previous section's \(V=V^{+}+V^{-}\) and \(Z_0I=V^{+}-V^{-}\) and the
two definitions collapse:
\[
a=\frac{V^{+}}{\sqrt{Z_0}},
\qquad
b=\frac{V^{-}}{\sqrt{Z_0}} .
\]
So \(a\) and \(b\) are nothing but the incident and reflected wave amplitudes,
rescaled. The rescaling is chosen so that the squared magnitudes are \emph{powers}:
with RMS phasors, \(|a|^2=|V^{+}|^2/Z_0\) is the power flowing in and \(|b|^2\) the
power coming back.

\begin{controlsbox}
\textbf{A change of basis, and an invertible one.} Solving backwards,
\(V=\sqrt{Z_0}\,(a+b)\) and \(I=(a-b)/\sqrt{Z_0}\), so
\[
\begin{bmatrix}a\\b\end{bmatrix}
=\frac{1}{2\sqrt{Z_0}}
\begin{bmatrix}1 & Z_0\\ 1 & -Z_0\end{bmatrix}
\begin{bmatrix}V\\I\end{bmatrix},
\qquad
\det=-\tfrac12\ne0 .
\]
No information is created or lost; \((a,b)\) and \((V,I)\) describe the same state in
two coordinate systems, exactly as \cref{sec:seriesparallel} rewrote one impedance two
ways and \cref{ch:rlc} found the poles were the eigenvalues whichever description you
started from. The reason to prefer wave coordinates at RF is practical: they are what
an instrument can actually separate, since a directional coupler responds to
\emph{direction of travel}, not to \(V\) or \(I\) individually.
\end{controlsbox}

For a one-port the whole story is one number, and it is one we already have:
\[
S_{11}=\frac{b}{a}=\frac{V^{-}}{V^{+}}=\Gamma
\]
\emph{identically}---not approximately, and not a new quantity. So \(|S_{11}|^2\) is
the reflected power fraction, return loss is \(-20\log_{10}|S_{11}|\), and SWR follows
from \(|S_{11}|\) by the formula of the next section. Three names for one measurement.

A two-port needs four numbers, because either port can be driven and either can
respond. Drive port 1 with port 2 terminated in \(Z_0\) (so nothing comes back into
port 2, \(a_2=0\)) and the two useful ones are
\[
S_{11}=\frac{b_1}{a_1}\bigg|_{a_2=0}
\quad\text{(what reflects)},
\qquad
S_{21}=\frac{b_2}{a_1}\bigg|_{a_2=0}
\quad\text{(what gets through)} .
\]
The subscripts are just \emph{out, in}: \(S_{21}\) is the wave out of port 2 per wave
into port 1. And \(S_{21}\) is the payoff, because it is a transfer function---the
forward transmission of the network as it sits in its \(Z_0\) environment. Everything
Chs \ref{ch:linsys} and \ref{ch:splane} said about \(G(\jj\omega)\) therefore applies
to it: magnitude and phase versus frequency, poles and zeros, decibels per decade,
group delay. \emph{A sweep of \(S_{21}\) is a Bode plot}, and a filter's
\(|S_{21}|\) trace is the response \cref{ch:filters} computes, measured instead. Its
in-band value is that chapter's \emph{insertion loss}
(\cref{sec:filterspecs}): \(-20\log_{10}|S_{21}|\), since \(|S_{21}|^2\) is the
fraction of incident power delivered. \Cref{sec:vna} is the instrument that performs
this sweep, and what it puts on its screen.

\begin{mistakebox}
\textbf{\(S_{21}\) is a wave ratio, not \(V_2/V_1\).} They coincide only when the
input is matched. In general \(V_1=\sqrt{Z_0}(a_1+b_1)\) includes the reflected wave,
so
\[
\frac{V_2}{V_1}=\frac{S_{21}}{1+S_{11}} ,
\]
and the two agree exactly when \(S_{11}=0\). For a filter doing its job in a
\SI{50}{\ohm} system that is the passband condition, which is why the identification
of \(S_{21}\) with the voltage transfer function is safe where you care and wrong at
the band edges. What never needs the caveat is power: \(|S_{21}|^2\) is the
transmitted power fraction whatever the match.
\end{mistakebox}

\begin{advancedbox}
\textbf{Why three calibration standards, and why those three.} A network analyzer
cannot measure \(a\) and \(b\) at the plane you care about; it measures them past its
own couplers, cables and connectors. All of that imperfection is a linear two-port
between the instrument and the device, and pushing an unknown \(\Gamma\) through a
linear two-port gives a bilinear relation
\[
\Gamma_{\text{meas}}=e_{00}+\frac{e_{\text{t}}\,\Gamma}{1-e_{11}\Gamma}
\]
with exactly \emph{three} independent complex unknowns---conventionally directivity
\(e_{00}\), source match \(e_{11}\), and reflection tracking \(e_{\text{t}}\). Three
unknowns need three equations, so three known standards suffice, and the natural three
are the ones the previous section already evaluated: a short (\(\Gamma=-1\)), a
matched load (\(\Gamma=0\)) and an open (\(\Gamma=+1\)). Measure those, solve the
three equations once, and every later measurement can be corrected back to the
reference plane. Calibration is not a tuning ritual; it is solving a \(3\times3\)
system whose right-hand side you obtained by measuring things you already knew.
\end{advancedbox}

\section{Standing Waves, SWR, and Return Loss}
Incident and reflected waves overlap everywhere along the line, and because they
travel in opposite directions their relative phase depends on position. At distance
\(x\) back from the load the backward wave has made a round trip of \(2x\), so
\[
|V(x)|=\bigl|V^{+}\bigr|\,\bigl|1+\Gamma\,\ee^{-\jj 2\beta x}\bigr| ,
\]
where \(\beta=2\pi/\lambda\) counts radians per unit length. The bracket sweeps
around a circle as \(x\) grows: where the two waves arrive in phase the magnitudes
\emph{add}, and where they arrive out of phase they \emph{subtract}. The pattern is
therefore stationary in space---a \textbf{standing wave}---with
\[
V_{\max}=|V^{+}|\bigl(1+|\Gamma|\bigr),
\qquad
V_{\min}=|V^{+}|\bigl(1-|\Gamma|\bigr),
\]
one quarter wavelength apart. Their ratio is the \textbf{standing-wave ratio},
and \(|V^{+}|\) cancels:
\[
\boxed{\;\mathrm{SWR}=\frac{V_{\max}}{V_{\min}}
=\frac{1+|\Gamma|}{1-|\Gamma|}\;}
\]
SWR equals 1 for a perfect match (\(\Gamma=0\) leaves no pattern at all) and rises
with mismatch, becoming infinite for a short or open. \textbf{Return loss} (RL)
expresses the same mismatch in the decibels of \cref{sec:decibels},
\[
\mathrm{RL}=-20\log_{10}|\Gamma| ,
\]
with higher return loss meaning a \emph{smaller} reflection. Because \(\Gamma\) is
an amplitude ratio it takes the \(20\log\) form; the reflected \emph{power}
fraction is \(|\Gamma|^2\).

\begin{mistakebox}
The symbol \(\mathrm{RL}\) for return loss has nothing to do with the RL circuits of
\cref{ch:rl}; the collision is unfortunate but standard. Likewise the phase constant
\(\beta\) in this chapter is unrelated to the feedback fraction \(\beta\) of
\cref{ch:feedback}---both notations are entrenched in their own fields, and only
context distinguishes them.
\end{mistakebox}

\section{Real Lines: Attenuation and the Complex Propagation Constant}
\label{sec:lineloss}
Everything so far has assumed a \emph{lossless} line, and that assumption has been
doing real work: it is why \(|\Gamma|\) was the same everywhere, why a stub was
purely reactive, and why SWR was a property of the load alone. Real coax
dissipates, so this section puts the loss back. It costs two more per-unit-length
parameters and one more derivation, and it repays them by explaining feed-line
loss, why that loss rises with frequency, and why a mismatch looks better at the
shack than it is at the antenna.

\subsection{Two More Parameters, One More Derivation}
The lossless slice had a series inductance \(L'\dd x\) and a shunt capacitance
\(C'\dd x\). A real slice adds the two dissipative elements alongside them: a
series resistance \(R'\dd x\) for the conductors, and a shunt conductance
\(G'\dd x\) for leakage through the dielectric. Applying KVL and KCL to the slice
as before---but now in phasor form, so \(\partial/\partial t\to\jj\omega\)---gives
the telegrapher's equations with loss:
\[
\frac{\dd V}{\dd x}=-\bigl(R'+\jj\omega L'\bigr)I,
\qquad
\frac{\dd I}{\dd x}=-\bigl(G'+\jj\omega C'\bigr)V .
\]
The structure is unchanged; each element law simply picked up its lossy companion.
Now run the same manoeuvre as before---differentiate the first with respect to
\(x\) and substitute the second:
\[
\frac{\dd^2V}{\dd x^2}
=-\bigl(R'+\jj\omega L'\bigr)\frac{\dd I}{\dd x}
=\underbrace{\bigl(R'+\jj\omega L'\bigr)\bigl(G'+\jj\omega C'\bigr)}_{\gamma^2}V .
\]
So \(V(x)\) obeys \(\dd^2V/\dd x^2=\gamma^2V\), whose solutions are
\(\ee^{\mp\gamma x}\), and the constant that governs everything is the
\textbf{complex propagation constant}
\[
\boxed{\;
\gamma=\sqrt{\bigl(R'+\jj\omega L'\bigr)\bigl(G'+\jj\omega C'\bigr)}
=\alpha+\jj\beta \;}
\]
Splitting \(\gamma\) into its real and imaginary parts is what makes it useful,
because a forward wave is then
\[
V(x)=V^{+}\ee^{-\gamma x}
=V^{+}\underbrace{\ee^{-\alpha x}}_{\text{decays}}
\underbrace{\ee^{-\jj\beta x}}_{\text{rotates}} .
\]
The two parts do independent jobs: \(\beta\) is the \emph{phase constant} already
met in the standing-wave section above, counting radians of delay per unit length, while
\(\alpha\) is the \textbf{attenuation constant}, shrinking the amplitude
exponentially with distance. Loss and delay, one complex number.

The same two parameters also make the characteristic impedance complex:
\[
Z_0=\sqrt{\frac{R'+\jj\omega L'}{G'+\jj\omega C'}} .
\]
Setting \(R'=G'=0\) recovers everything from before---\(\gamma=\jj\omega\sqrt{L'C'}\),
so \(\alpha=0\) and \(\beta=\omega\sqrt{L'C'}\), and \(Z_0=\sqrt{L'/C'}\) real. The
lossless line was the special case all along.

\subsection{The Low-Loss Limit, Which Is the Practical One}
Good coax is close to lossless: \(R'\ll\omega L'\) and \(G'\ll\omega C'\) across its
useful range. Expanding \(\gamma\) to first order in the small quantities---using
\(\sqrt{1+u}\approx1+u/2\) twice---gives
\[
\gamma\approx\jj\omega\sqrt{L'C'}
\left(1+\frac{R'}{2\jj\omega L'}+\frac{G'}{2\jj\omega C'}\right),
\]
and separating real from imaginary parts leaves the two results worth remembering:
\[
\boxed{\;
\alpha\approx\frac{R'}{2Z_0}+\frac{G'Z_0}{2},
\qquad
\beta\approx\omega\sqrt{L'C'},
\qquad
Z_0\approx\sqrt{\frac{L'}{C'}} \;}
\]
So in the low-loss regime the phase behavior and the impedance are exactly the
lossless ones---which retroactively justifies every earlier result---and the loss
appears only in \(\alpha\), as two independent contributions: a \emph{conductor}
term \(R'/2Z_0\) and a \emph{dielectric} term \(G'Z_0/2\).

\begin{physicalbox}
\textbf{Why loss rises with frequency, and at two different rates.} The two terms in
\(\alpha\) have different frequency dependence, and both come from physics already
met.

Conductor loss grows because of \emph{skin effect}: at higher frequency the current
crowds into a thinner outer layer of the conductor, so the usable cross-section
shrinks and \(R'\) rises as \(\sqrt f\). That exponent is not a fit---
\cref{sec:selfresonance} derived it from the skin depth
\(\delta=\sqrt{2/\omega\mu\sigma}\), and a conductor carrying its current in a layer
of thickness \(\delta\) has \(R\propto1/\delta\propto\sqrt f\). Dielectric loss grows because the
insulation's leakage tracks the displacement current through it, making \(G'\)
roughly proportional to \(f\).

Hence coax attenuation goes as \(\sqrt f\) where conductor loss dominates---the
entire HF and VHF range for ordinary cable---and crosses over to a steeper,
proportional-to-\(f\) behavior at microwave frequencies, where the dielectric term
takes over. That is the whole explanation of the shape of a manufacturer's
attenuation-versus-frequency curve, and \cref{fig:line-attenuation} separates the
two contributions.
\end{physicalbox}

\subsection{Attenuation in Decibels}
Manufacturers do not quote \(\alpha\) in nepers per metre; they quote loss in
decibels per unit length, conventionally per \SI{100}{\foot} or per
\SI{100}{\meter}. The conversion is just \cref{sec:decibels} applied to
\(\ee^{-\alpha x}\):
\[
\text{loss in dB}=-20\log_{10}\ee^{-\alpha x}=20\,\alpha x\log_{10}\ee
\approx8.686\,\alpha x .
\]
Because it is a decibel figure it is \emph{proportional to length}, which is what
makes the per-100-foot quotation useful: double the run, double the decibels. It also
costs more than it looks: by \cref{sec:noisefigure}, a passive loss ahead of a
receiver adds its full value to the system noise figure, so \SI{2}{\decibel} of feed
line is \SI{2}{\decibel} of sensitivity thrown away---which is the argument for
mounting a preamplifier at the antenna rather than in the shack
(\cref{sec:friis}).
Representative figures for common cable, which vary between manufacturers and
should be read off the actual data sheet:

\begin{center}
\begin{tabularx}{\textwidth}{@{}L{0.22\textwidth}Y Y Y@{}}
\toprule
Cable & Matched loss per \SI{100}{\foot} at \SI{14}{\mega\hertz} & \dots\ and at \SI{144}{\mega\hertz} & \(Z_0\) \\
\midrule
RG-58 (small coax) & \(\approx\SI{1.5}{\decibel}\) & \(\approx\SI{5}{\decibel}\) & \SI{50}{\ohm} \\
RG-213 (large coax) & \(\approx\SI{0.6}{\decibel}\) & \(\approx\SI{2}{\decibel}\) & \SI{50}{\ohm} \\
Open-wire / ladder & \(\approx\SI{0.1}{\decibel}\) & \(\approx\SI{0.4}{\decibel}\) & \SI{450}{\ohm} \\
\bottomrule
\end{tabularx}
\end{center}

\begin{figure}[htbp]
\centering
\includegraphics[width=\textwidth]{line_attenuation}
\caption{Feed-line attenuation. \emph{Left:} the two terms of
\(\alpha\approx R'/(2Z_0)+G'Z_0/2\) plotted separately for RG-58. On log--log
axes their different physics shows as different \emph{slopes}---the conductor
term rises as \(\sqrt f\) (slope \(\tfrac12\)) because of skin effect, the
dielectric term as \(f\) (slope \(1\)). Across the shaded amateur range the
conductor term is essentially the whole loss; the dielectric term is steeper and
takes over far above it. \emph{Right:} the same two-term model for both cables,
with the four representative values from the table marked. Computed from
\(\text{loss}=k_c\sqrt{f}+k_df\) with \((k_c,k_d)\) fitted to those values;
illustrative, not a substitute for a data sheet.}
\label{fig:line-attenuation}
\end{figure}

\begin{workedbox}[: does the \(\sqrt f\) law hold?]
Test the physics against the table. Between \SI{14}{\mega\hertz} and
\SI{144}{\mega\hertz} the frequency ratio is \(144/14\approx10.3\), so if conductor
loss dominates the attenuation should rise by
\[
\sqrt{10.3}\approx3.2\times .
\]
For RG-58 the table gives \(5/1.5\approx3.3\times\); for RG-213,
\(2.0/0.6\approx3.3\times\). Both match the \(\sqrt f\) prediction closely, which
confirms that ordinary coax is conductor-loss dominated through VHF---and that the
dielectric term has not yet begun to matter. Note also why open-wire line is so
much better: almost all of its field is in air, so \(G'\) is tiny, and its high
\(Z_0\) shrinks the conductor term \(R'/2Z_0\) as well.
\end{workedbox}

\subsection{Loss on a Mismatched Line}
Loss and mismatch interact, and in two directions worth separating.

\emph{Extra loss.} On a matched line a single wave travels once to the load. On a
mismatched line the reflected wave travels the length again, and the standing-wave
pattern concentrates current at some points and voltage at others---and since the
loss mechanisms are \(I^2R'\) and \(V^2G'\), those peaks dissipate more than a flat
distribution would. The total loss therefore exceeds the matched-line loss. Writing
\(a=10^{ML/10}\) for the matched-line loss \(ML\) expressed as a power ratio, the
total is
\[
\text{total loss (dB)}=10\log_{10}
\frac{a^2-|\Gamma_L|^2}{a\bigl(1-|\Gamma_L|^2\bigr)} .
\]
The formula is easier to trust after two checks. With \(|\Gamma_L|=0\) it collapses
to \(10\log_{10}(a^2/a)=ML\), the matched-line loss, as it must. And with \(a=1\) (a
\emph{lossless} line) it gives \(10\log_{10}1=0\) for \emph{any} \(|\Gamma_L|\)---so
high SWR on a lossless line costs nothing at all. That second limit is the important
one conceptually: SWR does not itself consume power. It only aggravates loss that the
line was already going to have, which is why a high SWR is a minor matter on short
runs of good cable at HF and a serious one on long runs at VHF.

\emph{Masked mismatch.} The reflected wave is also attenuated on its way back, so
the reflection coefficient seen at the input is smaller than the one at the load:
\[
\Gamma_{\mathrm{in}}=\Gamma_L\,\ee^{-2\gamma\ell}
\qquad\Longrightarrow\qquad
|\Gamma_{\mathrm{in}}|=|\Gamma_L|\,\ee^{-2\alpha\ell} .
\]
The factor of two is the round trip, so a line with \SI{3}{\decibel} of matched-line
loss attenuates the reflected wave by \SI{6}{\decibel}---halving \(|\Gamma|\), since
\(20\log_{10}2\approx\SI{6}{\decibel}\). \Cref{ch:measurement} works that arithmetic
through and draws the practical conclusion: an SWR meter at the transmitter can
report a comfortable match while the antenna is badly mismatched.

\begin{mistakebox}
Two statements about SWR are both true and easily confused. On a \emph{lossless}
line the SWR is the same everywhere, because \(|\Gamma|\) does not change---only its
phase rotates. On a \emph{lossy} line the SWR falls as you move away from the load,
because \(|\Gamma|\) itself decays as \(\ee^{-2\alpha\ell}\). The first is the
idealization the Smith chart's constant-SWR circles assume; the second is what a
real meter reads.
\end{mistakebox}

\section{Input Impedance Along the Line}
Because the reflected wave changes phase as it travels, the impedance seen
looking into a lossless line of length \(\ell\) toward a load \(Z_L\) rotates with
position:
\[
Z_{\mathrm{in}}(\ell)=Z_0\,
\frac{Z_L+\jj Z_0\tan\beta\ell}{Z_0+\jj Z_L\tan\beta\ell},
\qquad
\beta=\frac{2\pi}{\lambda}.
\]
The phase constant \(\beta\) counts radians of electrical length per unit
distance. The familiar special cases are limits of this one formula:
\begin{description}
\item[Quarter-wave (\(\beta\ell=\pi/2\), \(\tan\to\infty\))]
\(Z_{\mathrm{in}}=Z_0^2/Z_L\). The line \emph{inverts} the load impedance, which
is why \(Z_t=\sqrt{Z_0 Z_L}\) matches two real resistances.
\item[Half-wave (\(\beta\ell=\pi\), \(\tan=0\))] \(Z_{\mathrm{in}}=Z_L\): the line
repeats the load impedance.
\item[Matched (\(Z_L=Z_0\))] \(Z_{\mathrm{in}}=Z_0\) for any length: nothing to
rotate.
\end{description}
\section{Stubs: Line Sections as Reactances---and as Resonators}
\label{sec:stubs}
A \textbf{stub} is a short piece of line, left deliberately unterminated, used as a
circuit element rather than as a signal path. Its two useful cases are the extremes
of the formula above.

Setting \(Z_L=0\) (a \emph{shorted} stub) makes the numerator \(\jj Z_0\tan\beta\ell\)
and the denominator \(Z_0\), so
\[
Z_{\mathrm{in}}^{\text{short}}=\jj Z_0\tan\beta\ell ,
\]
and letting \(Z_L\to\infty\) (an \emph{open} stub) leaves the reciprocal behavior,
\[
Z_{\mathrm{in}}^{\text{open}}=-\jj Z_0\cot\beta\ell .
\]
Both are \emph{purely reactive}: a lossless stub cannot dissipate, so it can only
store. This makes a length of line a substitute for a capacitor or an inductor---
often a better one at VHF and above, where a real coil's parasitics dominate.

Reading the tangent tells you which. A shorted stub shorter than \(\lambda/4\) has
\(\tan\beta\ell>0\) and so looks \textbf{inductive}; an open stub shorter than
\(\lambda/4\) looks \textbf{capacitive}. Each behavior inverts every quarter
wavelength thereafter.

One length deserves singling out. At \(\ell=\lambda/8\) we have
\(\beta\ell=45^\circ\), where \(\tan=\cot=1\), so
\[
|X|=Z_0
\qquad\text{at}\qquad
\ell=\frac{\lambda}{8}
\]
for \emph{either} kind of stub---inductive if shorted, capacitive if open. And it is
the \emph{only} length below \(\lambda/4\) that does this, because \(\tan\) is
strictly increasing on \((0^\circ,90^\circ)\) and can equal \(1\) just once. That
makes \(\lambda/8\) the natural unit of stub design: it is where a stub's reactance
magnitude matches the line it is built from, so a \(\lambda/8\) shorted stub in a
\SI{50}{\ohm} system is a \(+\jj\SI{50}{\ohm}\) inductor with no calculation at all.

\begin{physicalbox}
\textbf{A stub is the two resonators of \cref{ch:rlc,ch:rlcpar} in disguise.}
Watch what happens at the special lengths:
\begin{itemize}
\item A \textbf{quarter-wave shorted} stub has \(\beta\ell=\pi/2\), where
\(\tan\to\infty\), so \(Z_{\mathrm{in}}\to\infty\). A short at the far end appears
as an \emph{open} at the near end---exactly what the parallel tank of
\cref{ch:rlcpar} does at \(\omega_0\). Near that frequency the stub \emph{is} a
parallel resonator.
\item A \textbf{half-wave shorted} stub has \(\beta\ell=\pi\), where \(\tan=0\), so
\(Z_{\mathrm{in}}=0\): a low impedance, which is what the series resonator of
\cref{ch:rlc} presents at \(\omega_0\).
\end{itemize}
And because \(\tan\beta\ell\) blows up not once but at every
\(\beta\ell=\pi/2+n\pi\), a single piece of line is an \emph{infinite ladder} of
alternating parallel and series resonances, spaced every quarter wavelength in
electrical length---equivalently, every odd and even multiple of a base frequency.
That is the concrete meaning of the ``infinite set of natural modes'' promised in
this chapter's opening: a lumped RLC circuit has one pole pair, while a line has
infinitely many, and the tangent is how they show up in a closed-form expression.
\Cref{sec:puredelay} anticipated exactly this from the other direction: a pure
delay has unbounded phase lag at undiminished magnitude, so it must cross every
multiple of \(-180^\circ\), and a line is a delay. \Cref{sec:infinitepoles} shows
the two pictures meeting---a cascade of \(n\) identical lags becomes a pure delay
as \(n\to\infty\), which is the lumped ladder of this section taken to its limit.

This is also why a stub makes a trap. \Cref{ch:filters} built notches from lumped
LC branches; a quarter-wave shorted stub bridged across a feed line does the same
job with no components at all, presenting an open at the design frequency and a low
impedance at its even harmonics.
\end{physicalbox}

\section{Quarter-Wave and Half-Wave Transformers}
\label{sec:quarterwave}
The quarter-wave limit is worth carrying one step further, because it is the
transformer formula the exam asks for. At \(\beta\ell=\pi/2\) the line inverts the
load about \(Z_0^2\),
\[
Z_{\mathrm{in}}=\frac{Z_0^2}{Z_L}.
\]
We want that input impedance to equal the system impedance we are feeding. Call the
\emph{section's} own characteristic impedance \(Z_t\) (it differs from the main
line's, which is the whole point) and require \(Z_{\mathrm{in}}=Z_0\):
\[
\frac{Z_t^2}{Z_L}=Z_0
\qquad\Longrightarrow\qquad
\boxed{\;Z_t=\sqrt{Z_0Z_L}\;}
\]
so the matching section's impedance is the \emph{geometric mean} of the two
resistances it joins. A half-wave section, by contrast, repeats the load
(\(\tan=0\)) and transforms nothing---useful for relocating a load electrically
without changing it.

Both results assume a lossless line and hold only near the design frequency, since
\(\beta\ell\) is frequency dependent---a narrowband match, as
\cref{fig:line-worked} makes quantitative.

\section{Smith Chart Interpretation}
A Smith chart is a graphical map of normalized complex impedance and reflection
coefficient. It plots the bilinear (M\"obius) transform between normalized
impedance \(z=Z/Z_0\) and \(\Gamma=(z-1)/(z+1)\), which turns resistance and
reactance grid lines into circles and arcs inside the unit \(|\Gamma|\le1\) disk.
Crucially, the \(Z_{\mathrm{in}}(\ell)\) rotation above becomes a simple rotation
about the chart's center: moving along the line spins the point on a
constant-\(|\Gamma|\) (constant-SWR) circle. For this book the point is
conceptual---the chart is not a new law of circuits but a coordinate transform
that makes the input-impedance formula something you can read off by eye.

\Cref{fig:smith-labelled} is the chart with every feature the exam names labelled, and
the vocabulary is worth learning deliberately because two of the names are
counter-intuitive.

Start with the \emph{two} families, which are all the chart is: circles of constant
\textbf{resistance} and arcs of constant \textbf{reactance}. Nothing else is plotted.
The resistance circles are all tangent to the right-hand point, where \(z\to\infty\);
the reactance arcs all terminate on the outer circle.

Now the two names that trip people. The horizontal line through the middle is the
\textbf{resistance axis}---and it is the \emph{only straight line on the chart}, which
is how a question can ask for it without naming it. But the large outer circle is
called the \textbf{reactance axis}, because that is where the reactance arcs
terminate---an ``axis'' that is a circle, which is exactly the sort of thing worth
seeing once on a labelled diagram rather than meeting first in a question. The center
is the \textbf{prime center}: \(z=1\), \(\Gamma=0\), the perfect match.

\emph{Normalizing} the chart is the operation that makes it usable in any system. The
chart is drawn once in normalized coordinates \(z=Z/Z_0\), so to work in a
\SI{50}{\ohm} system you declare the prime center to \emph{be} \SI{50}{\ohm}---you
reassign the prime center's impedance value---and every number you read off is then a
multiple of it. Change to \SI{75}{\ohm} and not one line moves; only the label on the
center does.

Two additions turn the chart from a picture into a tool. \textbf{Wavelength scales}
around the rim are calibrated in fractions of a \emph{transmission-line} wavelength,
not an antenna's, because what the rotation tracks is distance along the line: one full
turn is \(\lambda/2\), so a half turn is the \(\beta\ell=\pi/2\) that inverts an
impedance. And a third family of \textbf{constant-SWR circles}---concentric about the
prime center, since \(|\Gamma|\) and SWR determine each other---is what gets drawn on
top when the chart is used to design a matching network.

\subsection{What the Chart Is Actually For}
\label{sec:smithuse}
The chart's characteristic job is the one this chapter has already done twice
algebraically: \emph{find the length and position of a matching stub}. The recipe is
the rotation. Moving from the load toward the generator rotates the point clockwise on
its constant-SWR circle, so you rotate until the point lands on the circle where the
normalized \emph{resistance} is \(1\)---the distance rotated is where the stub goes,
read off the wavelength scale. Whatever reactance remains at that point is what the
stub must cancel, and \cref{sec:stubs}'s \(\jj Z_0\tan\beta\ell\) converts that
reactance into a stub length. Two readings off the chart, and the answer is the same
pair of numbers \cref{sec:antennamatches}'s stub match arrives at by algebra.

That is the honest summary of what a Smith chart is: not a new law, and not a
calculating aid anyone needs now that \(Z_{\mathrm{in}}(\ell)\) can be evaluated on a
phone, but a \emph{picture} in which the algebra's structure is visible---rotation for
distance, circles for constant SWR, one point for a match.

\begin{exambox}
Group E9G is eleven questions of vocabulary and use, and every answer is either
labelled in \cref{fig:smith-labelled} or stated in this section. What it \emph{calculates}:
impedance along transmission lines, and impedance and SWR values. Its \emph{coordinate
system}: resistance circles and reactance arcs---the two families, and the arcs are the
points of constant \emph{reactance}. Its \emph{common use}: determining the length and
position of a matching stub. The \emph{only straight line}: the resistance axis. The
\emph{large outer circle on which the reactance arcs terminate}: the reactance axis.
\emph{Normalizing}: reassigning the prime center's impedance value. The \emph{third
family added when designing matching networks}: constant-SWR circles. And the
\emph{wavelength scales}: fractions of a transmission-line electrical wavelength---the
distractors offer antenna wavelength and frequency, and the scale tracks distance along
the \emph{line}.
\end{exambox}

\begin{figure}[htbp]
\centering
\includegraphics[width=0.92\textwidth]{smith_chart_labelled}
\caption{A Smith chart with every feature the exam names labelled. It is not traced:
each curve is computed from \(\Gamma=(z-1)/(z+1)\), the constant-resistance circles
having centre \(r/(1+r)\) and radius \(1/(1+r)\) and the constant-reactance arcs centre
\((1,1/x)\) and radius \(|1/x|\). Note the two names worth learning: the
\emph{resistance} axis is the horizontal line---the only straight line on the
chart---while the \emph{reactance} axis is the outer \emph{circle}, where every arc
terminates. The orientation, with \(z\to\infty\) at the right, is the conventional one
the question pool adopted in its 2024 errata.}
\label{fig:smith-labelled}
\end{figure}


\section{An Idealized Antenna on the Smith Chart}
The most useful way to connect the chart to the math is to put a real device on
it. An antenna is where the RF energy finally leaves the circuit, and near its
fundamental resonance an antenna behaves like a one-port whose impedance we can
model with ideal elements---exactly the schematic~\(\to\) model~\(\to\) plot
pattern of the earlier circuit chapters.

\subsection{Modeling the Antenna}
Near resonance, a center-fed dipole looks electrically like a \textbf{series RLC}
one-port (\cref{fig:antenna-model}): a \emph{radiation resistance} \(R_{\text{rad}}\)
in series with an inductance and capacitance that together set the resonance. The
radiation resistance is not a lossy resistor---it represents power that leaves as
radiation---but electrically it behaves like one. That means the whole of
\cref{ch:rlc} applies directly: the impedance is that chapter's series-RLC impedance,
\[
Z(\omega)=R_{\text{rad}}+\jj\left(\omega L-\frac{1}{\omega C}\right),
\qquad
\omega_0=\frac{1}{\sqrt{LC}} ,
\]
its quality factor is that chapter's \(Q_s=\omega_0L/R\), and its half-power
bandwidth is \(BW=f_0/Q\). \Cref{sec:halfpower} showed that equality to be exact,
so the \(\approx\) an antenna's bandwidth deserves belongs to the \emph{model}: a
real antenna's \(R_{\text{rad}}\), \(L\) and \(C\) drift with frequency, whereas the
lumped RLC holds them fixed.

Below \(\omega_0\) the antenna looks capacitive (\(X<0\), the ``too short''
regime); above it, inductive (\(X>0\), ``too long''); at \(\omega_0\) the
reactance vanishes and \(Z=R_{\text{rad}}\) is purely real. A half-wave dipole is
``resonant'' precisely because it is the length at which those two cancel.

\subsection{From Wavelength to Wire Length}
\label{sec:antennalength}
That last sentence has a practical corollary: the resonant length is a
\emph{wavelength} statement, so it can be turned into a number with a tape measure.
The familiar formulas are worth deriving, because doing so shows where their odd
constants come from.

In free space \(\lambda=c/f\). Working in the units the formulas are quoted
in---feet and megahertz---and using
\(c\approx\SI{9.836e8}{\foot\per\second}\),
\[
\lambda\,[\mathrm{ft}]\approx\frac{984}{f\,[\mathrm{MHz}]},
\qquad\text{so}\qquad
\frac{\lambda}{2}\approx\frac{492}{f} .
\]
A real dipole is slightly \emph{shorter} than that. The wire has thickness and its
ends have capacitance to their surroundings---the \textbf{end effect}---which adds a
little capacitive loading, so resonance arrives at a physical length about
\SI{5}{\percent} short of the free-space half wavelength. Applying that factor:
\[
\boxed{\;
\ell_{\lambda/2}\,[\mathrm{ft}]\approx\frac{0.95\times492}{f}
\approx\frac{468}{f\,[\mathrm{MHz}]},
\qquad
\ell_{\lambda/4}\,[\mathrm{ft}]\approx\frac{234}{f\,[\mathrm{MHz}]} \;}
\]
the second being a quarter-wave monopole, half the dipole. So the \(468\) is not a
mysterious constant: it is \(984/2\) reduced by the end effect, and the \(234\) is
just half of it.

\begin{workedbox}[: cutting a dipole for 20 metres]
For \(f=\SI{14.1}{\mega\hertz}\), the antenna modeled above,
\[
\ell=\frac{468}{14.1}\approx\SI{33.2}{\foot}
\qquad\text{total, so about \SI{16.6}{\foot} per leg.}
\]
Compare the two \(\SI{5}{\percent}\) corrections that appear in this chapter and do
\emph{not} confuse them: the end-effect factor above applies to an antenna radiating
in space, while the velocity factor \(VF\) of \cref{sec:lineloss}'s cables applies to
a wave guided \emph{inside} a line. A \(\lambda/4\) matching \emph{stub} of RG-213 at
this frequency is \(234/14.1\times0.66\approx\SI{11}{\foot}\)---shortened by
\(VF\), not by end effect.
\end{workedbox}

\subsection{Gain, and Two Reference Antennas}
Antenna gain is a ratio, so \cref{sec:decibels} applies unchanged---but a ratio needs
a reference, and antenna work uses two:
\begin{description}
\item[dBi] gain relative to an \emph{isotropic} radiator, a hypothetical lossless
antenna radiating equally in every direction. It is a mathematical fiction, which is
what makes it a clean reference.
\item[dBd] gain relative to a half-wave \emph{dipole} in free space.
\end{description}
A dipole is not isotropic---it concentrates its radiation broadside, and that
concentration is worth \SI{2.15}{\decibel}. So the two scales differ by a constant:
\[
\boxed{\;\mathrm{dBi}=\mathrm{dBd}+2.15\;}
\]
which means a figure quoted in dBi always looks \SI{2.15}{\decibel} better than the
same antenna quoted in dBd. Nothing physical changes; only the reference does.

\begin{mistakebox}
Gain here is \emph{directivity}, not amplification. A passive antenna adds no energy:
\SI{10}{\decibel} of gain in one direction is bought by radiating less in others.
This is why gain and pattern are the same statement, and why an antenna cannot have
gain in every direction at once.
\end{mistakebox}

Those two references are all that separates the two ``radiated power'' figures a
regulation or a data sheet may quote. Both are the link budget of \cref{sec:sunits}
stopped after the antenna instead of carried to the far end---transmitter power, minus
what the feed line and connectors lose, plus what the antenna concentrates:
\[
\mathrm{ERP}\,[\dB\mathrm{W}]=P_{\text{tx}}-L_{\text{line}}+G_{\mathrm{dBd}},
\qquad
\mathrm{EIRP}\,[\dB\mathrm{W}]=P_{\text{tx}}-L_{\text{line}}+G_{\mathrm{dBi}} .
\]
The only difference is which reference the gain used, so
\(\mathrm{EIRP}=\mathrm{ERP}+\SI{2.15}{\decibel}\) always---a factor of \(1.64\) in
watts. A \SI{100}{\watt} transmitter through \SI{1.5}{\decibel} of feed line into a
\(6\,\mathrm{dBd}\) beam gives
\(\mathrm{ERP}=\SI{282}{\watt}\) and \(\mathrm{EIRP}=\SI{462}{\watt}\): the same
station, described twice. Note also that neither is efficiency: the \(\eta\) of
\cref{sec:antennaeff} below is folded into the gain figure, since a lossy antenna
concentrates less of what reaches it.

\subsection{Putting Numbers On It}
Take a dipole resonant in the \SI{20}{\meter} band at
\(f_0=\SI{14.1}{\mega\hertz}\), with radiation resistance
\(R_{\text{rad}}=\SI{35}{\ohm}\)---about half the \SI{73}{\ohm} of a dipole in
free space, which is what a real one at ordinary height presents, because the ground
image of \cref{sec:groundimage} couples back into it and pulls the feed resistance
down---and characteristic reactance
\(\omega_0L=1/(\omega_0C)=\SI{200}{\ohm}\). Then
\[
L=\frac{\SI{200}{\ohm}}{\omega_0}\approx\SI{2.26}{\micro\henry},
\qquad
C=\frac{1}{\omega_0\,(\SI{200}{\ohm})}\approx\SI{56}{\pico\farad},
\]
and---reusing \cref{ch:rlc}'s formulas unchanged---
\[
Q=\frac{\omega_0L}{R_{\text{rad}}}=\frac{200}{35}\approx5.7,
\qquad
BW=\frac{f_0}{Q}=\frac{\SI{14.1}{\mega\hertz}}{5.7}\approx\SI{2.5}{\mega\hertz}.
\]
That last number is the practical payoff of the whole model, and it answers a
question nothing earlier in this chapter could: \emph{what sets an antenna's usable
bandwidth?} Its \(Q\), exactly as for any other resonator. A \(Q\) near \(6\) is very
low by resonator standards---the pole pair sits far from the \(\jj\omega\) axis and
the resonance is broad---which is why a full-size dipole covers a whole amateur band.
Checking the band edges with \(\Gamma=(Z-Z_0)/(Z+Z_0)\) confirms it: at
\SI{14.0}{\mega\hertz}, \(Z\approx35-\jj2.8\ \si{\ohm}\) and
\(\mathrm{SWR}\approx1.4\); at \SI{14.35}{\mega\hertz},
\(Z\approx35+\jj7.0\ \si{\ohm}\) and \(\mathrm{SWR}\approx1.5\). The residual
\(\mathrm{SWR}\approx1.4\) even at resonance is not a reactance problem at all---it is
simply that \(35\,\si{\ohm}\ne50\,\si{\ohm}\), which is what a matching network from
\cref{ch:filters} would fix---or, in the forms antenna work gives proper names to,
one of the schemes of \cref{sec:antennamatches}.

\begin{physicalbox}
\textbf{Why a short loaded whip is narrowband.} Shorten the antenna and
\(R_{\text{rad}}\) collapses, while the loading coil needed to cancel the now-large
capacitive reactance pushes \(\omega_0L\) \emph{up}. Both changes drive
\(Q=\omega_0L/R_{\text{rad}}\) upward, so \(BW=f_0/Q\) shrinks---which is why a mobile
whip needs retuning across a band a dipole covers in one go, and why its efficiency
suffers (more of the loss now sits in the coil's resistance than in radiation). One
formula from \cref{ch:rlc}, two practical consequences.
\end{physicalbox}

\subsection{From the Model to the Chart}
Normalizing by the system impedance, \(z(\omega)=Z(\omega)/Z_0\), the reflection
coefficient is the same M\"obius map as before,
\[
\Gamma(\omega)=\frac{z(\omega)-1}{z(\omega)+1}.
\]
Sweeping frequency traces a curve on the Smith chart (\cref{fig:smith-antenna}).
At resonance the point crosses the real axis at
\(z=R_{\text{rad}}/Z_0\); if \(R_{\text{rad}}\ne Z_0\) it lands left of center
(here \(35/50=0.7\)), a residual mismatch a matching network would move to the
center. Off resonance the point swings up (inductive) or down (capacitive),
tracing the arc a NanoVNA draws as it sweeps the band. Reading the chart is
reading the model: distance from center is \(|\Gamma|\) (hence SWR), the crossing
of the real axis marks \(\omega_0\), and which half of the chart the arc sits in
tells you whether to add series capacitance or inductance to tune it.

\begin{controlsbox}
\textbf{One pole pair, three pictures.} These are worth lining up, because they are
the same series-RLC model drawn in three coordinate systems.
\begin{itemize}
\item In \cref{ch:rlc} the impedance locus \(Z(f)=R+\jj X(f)\) was a \emph{vertical
straight line} at \(\Re\{Z\}=R\) (\cref{fig:rlc-phasor}), because \(R\) is fixed and
only the reactance sweeps.
\item In \cref{ch:rlcpar} the dual circuit's locus was a \emph{circle through the
origin} (\cref{fig:rlcpar-worked}), because taking the reciprocal of a straight line
that misses the origin produces a circle that passes through it.
\item Here the Smith chart applies \(\Gamma=(z-1)/(z+1)\), and that map likewise
sends straight lines to circles---which is why the antenna sweep in
\cref{fig:smith-antenna} is an \emph{arc}.
\end{itemize}
Nothing physical differs among the three; only the coordinates do. It is the same
point \cref{ch:fourviews} made about time, state space, and frequency, now made about
three ways of plotting one impedance.
\end{controlsbox}

\begin{figure}[htbp]
\centering
\includegraphics[width=0.66\textwidth]{smith_chart_antenna}
\caption{Reflection-coefficient locus \(\Gamma(f)\) of the idealized series-RLC
antenna (\(R_{\text{rad}}=\SI{35}{\ohm}\), \(f_0=\SI{14.1}{\mega\hertz}\)) on a
hand-built Smith chart. The sweep crosses the real axis at resonance and swings
inductive (upper half) above \(f_0\), capacitive (lower half) below---the same
picture a NanoVNA shows for a real antenna.}
\label{fig:smith-antenna}
\end{figure}

\begin{exambox}
Almost every transmission-line question is one of four, and each has a single
governing formula. \emph{How bad is the mismatch?}
\(\Gamma=(Z_L-Z_0)/(Z_L+Z_0)\), then
\(\mathrm{SWR}=(1+|\Gamma|)/(1-|\Gamma|)\); for a purely resistive load this
shortcuts to the larger of \(Z_L/Z_0\) and \(Z_0/Z_L\). \emph{How long is the line
electrically?} \(\lambda=VF\,c/f\)---and use the real line's \(VF\), never
free-space wavelength. \emph{What does the line do to the impedance?} A quarter
wave inverts it about \(Z_0^2\), a half wave repeats it, a matched line leaves it
alone. \emph{What does an unterminated piece do?} It is a pure reactance: a shorted
stub under \(\lambda/4\) is inductive, an open stub under \(\lambda/4\) is
capacitive, and the behavior inverts every quarter wave.

Two traps. A perfect match is about \(Z_L=Z_0\), so no line length can fix a
mismatch---it only relocates it. And reflected \emph{power} is \(|\Gamma|^2\) while
return loss is \(-20\log_{10}|\Gamma|\), so watch whether a question is asking about
amplitude or power.
\end{exambox}


\section{What the Antenna Model Explains}
The same one-port model, pressed further, accounts for efficiency, feed position,
arrays, and the effect of ground---each by reusing a result this book already owns
rather than by introducing antenna theory of its own.

\subsection{Efficiency: the Resistance That Does Not Radiate}
\label{sec:antennaeff}
The model so far has one resistance, and it has been quietly doing two jobs. Split
them. An antenna's feed-point resistance is
\[
R=R_{\text{rad}}+R_{\text{loss}},
\]
where \(R_{\text{rad}}\) accounts for power that leaves as radiation and
\(R_{\text{loss}}\) for power that becomes heat---conductor resistance including the
skin effect of \cref{sec:selfresonance}, a loading coil's own loss, dielectric loss in
nearby insulators, and, for a ground-mounted vertical, the resistance of the return
currents flowing through soil. Both resistances carry the \emph{same} current, because
they are in series, so the power splits in proportion to them---\cref{ch:resistive}'s
divider, applied to power instead of voltage:
\[
\boxed{\;\eta=\frac{P_{\text{radiated}}}{P_{\text{input}}}
=\frac{I^2R_{\text{rad}}}{I^2\left(R_{\text{rad}}+R_{\text{loss}}\right)}
=\frac{R_{\text{rad}}}{R_{\text{rad}}+R_{\text{loss}}}\;}
\]
Nothing here is new physics; it is a divider. But it explains the two facts that
matter about small antennas. A full-size dipole has \(R_{\text{rad}}\approx\SI{73}{\ohm}\)
against a fraction of an ohm of loss, so \(\eta\) is above \SI{95}{\percent} and the
distinction hardly matters. Shorten the antenna and \(R_{\text{rad}}\) falls
roughly as the square of its length while \(R_{\text{loss}}\) does not, so the ratio
collapses. That is also why a ground radial system helps a quarter-wave vertical: the
radials do not change \(R_{\text{rad}}\), they shrink the \(R_{\text{loss}}\)
that soil return currents contribute.

\begin{mistakebox}
\textbf{Loss makes an antenna look \emph{better} on an SWR meter.} This is the
counterintuitive one, and both halves of it follow from formulas already derived. Take
a short antenna with \(R_{\text{rad}}=\SI{15}{\ohm}\) in a \SI{50}{\ohm} system and
add loss:

\begin{center}
\begin{tabular}{@{}rrrrr@{}}
\toprule
\(R_{\text{loss}}\) & \(R\) total & \(\eta\) & Signal lost & SWR at \(f_0\) \\
\midrule
\SI{0}{\ohm}  & \SI{15}{\ohm} & \SI{100}{\percent} & \SI{0}{\decibel}    & \(3.3{:}1\) \\
\SI{10}{\ohm} & \SI{25}{\ohm} & \SI{60}{\percent}  & \SI{-2.2}{\decibel} & \(2.0{:}1\) \\
\SI{35}{\ohm} & \SI{50}{\ohm} & \SI{30}{\percent}  & \SI{-5.2}{\decibel} & \(1.0{:}1\) \\
\bottomrule
\end{tabular}
\end{center}

At resonance the reactances cancel and the SWR is just \(Z_0/R\) or \(R/Z_0\), so
piling on loss walks \(R\) toward \SI{50}{\ohm} and the match toward perfection. The
bandwidth improves too, and exactly: \(Q=\omega_0L/R\) falls as \(R\) rises, and
\cref{sec:halfpower} made \(BW=f_0/Q\) an equality, so the usable window widens in
direct proportion to the loss added. A dummy load is the limiting case---a flawless
\(1{:}1\) across every band, radiating nothing. \emph{An SWR meter measures how
comfortable the transmitter is, not how much signal leaves.} That is the whole reason
\(\eta\) has to be a separate number, and why it cannot be measured from the shack.
\Cref{sec:travelingwave} is the case where that trade is made on purpose.
\end{mistakebox}

\begin{exambox}
E9A asks for efficiency as \emph{radiation resistance divided by total resistance},
which is the boxed formula, and separately asks what improves a ground-mounted
vertical: \emph{a ground radial system}, because it attacks \(R_{\text{loss}}\).
Watch the distractor that defines efficiency as ERP over transmitter output---that is
a different ratio, and it folds in feed-line loss and antenna gain rather than the
feed-point split.
\end{exambox}

\subsection{Where You Feed It}
\label{sec:feedpoint}
So far the antenna has been fed at its centre. Move the feed and the resistance
changes---predictably, and by an argument that needs no electromagnetics beyond what
\cref{sec:stubs} already set up.

A half-wave dipole is an open-ended structure, so the same boundary condition that put
a current \emph{node} at the end of an open stub puts one at each tip of the dipole.
The current is therefore a standing wave, largest at the centre and vanishing at the
ends:
\[
I(z)=I_0\cos\beta z,
\qquad |z|\le\frac{\lambda}{4},
\]
with \(z\) measured from the centre and \(\beta=2\pi/\lambda\) as always. Now the key
step, and it is conservation of power rather than a field calculation. The radiated
power is a property of that whole current distribution, so it does not care where you
attach the coax. But at the feed point the transmitter sees
\(P=|I(z)|^2R(z)\) with RMS values. One power, two ways of writing it:
\[
|I(z)|^2R(z)=\text{const}
\qquad\Longrightarrow\qquad
\boxed{\;R(z)=\frac{R_{\text{centre}}}{\cos^2\beta z}\;}
\]
Feed where the current is small and the same power demands a larger resistance. This
assumes moving the feed does not reshape the current distribution, which is the
standard approximation and a good one away from the extreme ends.

That one formula generates the whole subject:

\begin{center}
\begin{tabular}{@{}llrr@{}}
\toprule
Feed position & \(\beta z\) & \(\cos^2\beta z\) & \(R\) \\
\midrule
Centre & \(0^\circ\) & \(1\) & \SI{73}{\ohm} \\
\(\lambda/16\) from centre & \(22.5^\circ\) & \(0.854\) & \SI{85}{\ohm} \\
\(\lambda/8\) from centre & \(45^\circ\) & \(\tfrac12\) & \SI{146}{\ohm} \\
\(0.2\lambda\) from centre & \(72^\circ\) & \(0.095\) & \SI{765}{\ohm} \\
End & \(90^\circ\) & \(0\) & \(\to\infty\) \\
\bottomrule
\end{tabular}
\end{center}

The \(\lambda/8\) row is exact and worth remembering as a sanity check:
\(\cos^245^\circ=\tfrac12\), so the resistance is exactly \emph{doubled}, \SI{146}{\ohm}.
An off-centre-fed dipole works in this region, where the impedance is a few hundred
ohms rather than \SI{73}{\ohm}, so it is reached through a step-up balun rather than
directly. Its actual selling point is a multiband one: at the chosen offset the feed
point lands at a comparable fraction of the standing wave on several harmonically
related bands, giving a \emph{similar} impedance on each---which is a statement about
where the nodes fall, not about any one band's number.

The last row is the interesting one. \(R\to\infty\) at the end means an end-fed half
wave is a \emph{voltage} feed: there is almost no current to drive, and no ordinary
line can match it directly. But \cref{sec:quarterwave} already derived the cure. A
quarter-wave section inverts, \(Z_{\mathrm{in}}=Z_0^2/Z_L\), so it turns that very
high impedance into a very low one---which is why an end-fed half wave is fed through
a quarter-wave transformer, a high-ratio transformer, or an L-network doing the same
job. The mechanism is not special to antennas; it is the impedance inversion of
\cref{sec:quarterwave} used deliberately.

\begin{exambox}
E9A asks which factors affect feed-point impedance, and the answer is \emph{antenna
height}---not line length, not tuner settings, not power level. This section explains
why the first two are excluded: the feed-point impedance is a property of the current
distribution on the antenna, so it changes when the antenna's surroundings change the
distribution (height above ground does, via the ground reflection of
\cref{sec:groundimage}) and not when you
alter something on the transmitter side of the feed point. A line merely
\emph{transforms} what is already there (\cref{ch:lines}), which is a different claim
from changing it.
\end{exambox}

\subsection{Two Elements: the Array Factor}
\label{sec:arrayfactor}
Put two identical antennas a distance \(d\) apart and feed them with a deliberate
phase difference \(\phi\). The pattern that results is usually presented as a table of
four cases to memorize. It is one line of phasor addition.

Look in a far-off direction making angle \(\theta\) with the line joining the two
elements. The wave from the far element travels an extra distance \(d\cos\theta\) to
reach that direction, which is a phase lag \(\beta d\cos\theta\); add the feed phase
\(\phi\) you chose. Both elements radiate the same pattern, so that common factor
divides out and what remains is the sum of two unit phasors---\cref{ch:complex}'s
addition, nothing more:
\[
F(\theta)=1+\ee^{\jj\left(\phi+\beta d\cos\theta\right)} .
\]
Factor out half the angle, exactly as when adding any two phasors:
\(1+\ee^{\jj x}=\ee^{\jj x/2}\bigl(\ee^{-\jj x/2}+\ee^{\jj x/2}\bigr)=2\ee^{\jj x/2}\cos(x/2)\).
The leading exponential has unit magnitude, so
\[
\boxed{\;|F(\theta)|=2\left|\cos\!\left(\frac{\phi+\beta d\cos\theta}{2}\right)\right|\;}
\]
Every two-element pattern is now a substitution. With \(\beta d=2\pi d/\lambda\):

\begin{center}
\begin{tabularx}{\textwidth}{@{}llccY@{}}
\toprule
\(d\) & \(\phi\) & \(|F|\) along axis & \(|F|\) broadside & Pattern \\
\midrule
\(\lambda/2\) & \(0^\circ\)   & \(2|\cos90^\circ|=0\) & \(2|\cos0^\circ|=2\) & figure-eight \emph{broadside} to the axis \\
\(\lambda/2\) & \(180^\circ\) & \(2|\cos180^\circ|=2\) & \(2|\cos90^\circ|=0\) & figure-eight \emph{along} the axis (end-fire) \\
\(\lambda/4\) & \(90^\circ\)  & \(2|\cos90^\circ|=0\) & \(2|\cos45^\circ|=\sqrt2\) & \emph{cardioid}: one lobe, one null \\
\bottomrule
\end{tabularx}
\end{center}

Read the mechanism off the formula rather than the table. Feeding in phase makes the
two contributions add broadside, where the path difference is zero, and cancel along
the axis whenever the half-wave spacing turns the path difference into a half cycle.
Feeding out of phase swaps the two. And the cardioid case is the pretty one: with
\(d=\lambda/4\) the path difference along the axis is a quarter cycle, which the
\(90^\circ\) feed phase either doubles to a half cycle---complete cancellation on one
side---or annuls to zero, giving full addition on the other. \emph{One deliberate
phase and one path length, set to agree in one direction and disagree in the other.}
\Cref{fig:array-patterns} plots all three.

\begin{figure}[htbp]
\centering
\includegraphics[width=\textwidth]{array_patterns}
\caption{The three two-element patterns the pool names, all of them
\(|F|=2|\cos((\phi+\beta d\cos\theta)/2)|\) evaluated with different \(d\) and
\(\phi\)---no new physics between the panels, only two substitutions. Feeding in phase
at half-wave spacing cancels along the axis and adds broadside; reversing the phase
swaps the two; and quarter-wave spacing with a \(90^\circ\) feed phase cancels
completely on one side, giving the cardioid. The cardioid panel also carries the two
measurements E9B asks to read off a pattern---the \(-\SI{3}{\decibel}\) beamwidth and
the front-to-back null---taken from the same curve rather than a separate diagram. In
every panel the two elements are the grey squares, lying on the horizontal array
axis.}
\label{fig:array-patterns}
\end{figure}

\begin{physicalbox}
\textbf{Why direction finding listens for the null.} Near a maximum the cosine is
flat, so the signal barely changes as you rotate; near a zero it is steep, and the
signal falls off a cliff. A cardioid's single null is therefore a far sharper
indicator of bearing than its single broad lobe, which is why a direction-finding
set---a loop plus a sense antenna, phased to make exactly the third row above---is
aimed by nulling the signal rather than peaking it. The sharpness is not a property of
the hardware but of differentiating a cosine near its zero.
\end{physicalbox}

\begin{exambox}
Group E9C asks for three of these patterns by name, and all three are the table above:
two quarter-wave verticals \(\lambda/2\) apart \emph{in phase} give a figure-eight
\emph{broadside} to the array axis; the same pair fed \(180^\circ\) out of phase give
a figure-eight \emph{along} the axis; and \(\lambda/4\) apart fed \(90^\circ\) out of
phase give a \emph{cardioid}. Rather than memorizing which is which, evaluate
\(|F|=2|\cos((\phi+\beta d\cos\theta)/2)|\) at \(\theta=0^\circ\) and
\(90^\circ\)---two arithmetic steps, and the pattern name follows from where the zero
lands.
\end{exambox}

\subsection{Ground Is an Image, and Three Antennas That Follow}
\label{sec:groundimage}
\Cref{sec:arrayfactor} was written for two real antennas, but it does not care whether
the second one exists. A conducting ground plane forces the tangential electric field
to zero at its surface, and that boundary condition is satisfied exactly by adding an
\emph{image} antenna as far below the ground as the real one is above it. So an antenna
at height \(h\) over ground \emph{is} a two-element array with spacing \(d=2h\), and
the only question is the image's phase.

For a \textbf{horizontal} antenna the image is inverted, \(\phi=180^\circ\). Measuring
elevation \(\psi\) up from the horizon, the vertical separation projects as
\(2h\sin\psi\), so \cref{sec:arrayfactor} gives
\[
|F|=2\left|\cos\!\left(\frac{\pi+2\beta h\sin\psi}{2}\right)\right|
=2\left|\sin\left(\beta h\sin\psi\right)\right| .
\]
Two facts drop out. At \(\psi=0\) the sine vanishes: \emph{a horizontal antenna over
ground radiates nothing at the horizon}, at any height, because the image always
cancels it there. And the first maximum sits where \(\beta h\sin\psi=\pi/2\), that is
\[
\sin\psi=\frac{\lambda}{4h},
\]
so height alone sets the lobe: \(h=\lambda/4\) puts it straight up, \(h=\lambda/2\)
at \(30^\circ\), \(h=\lambda\) at \(14.5^\circ\), and \(h=1.5\lambda\) at
\(9.6^\circ\). Raising a horizontal antenna lowers its lobe---which is the whole
reason height matters for DX, and it is arithmetic, not folklore. \Cref{fig:ground-image}
plots both families.

For a \textbf{vertical} antenna the image is \emph{not} inverted, \(\phi=0\), and the
same formula returns a \emph{maximum} at the horizon instead. That single sign is why
verticals are favoured for low-angle work and why a horizontal antenna must be got up
high to compete. It is also why \cref{sec:feedpoint}'s answer to ``what affects
feed-point impedance'' was \emph{antenna height}: the image is a second element
coupled to the first, and moving it changes the impedance as surely as moving a real
one would.

\begin{figure}[htbp]
\centering
\includegraphics[width=\textwidth]{ground_image_lobes}
\caption{Ground as an image antenna, drawn with the array factor of
\cref{sec:arrayfactor}. \emph{Left:} a horizontal antenna, whose image is inverted, so
\(|F|=2|\sin(\beta h\sin\psi)|\). Every curve goes to zero at the horizon---the image
cancels the antenna there at any height---and the first lobe descends as the antenna
rises, from straight up at \(\lambda/4\) to \(30^\circ\) at \(\lambda/2\) and
\(14^\circ\) at \(\lambda\), following \(\sin\psi=\lambda/4h\). At \(h=\lambda\) a
second lobe has appeared. \emph{Right:} a vertical antenna, whose image is \emph{not}
inverted, so the sine becomes a cosine and the maximum lands \emph{at} the horizon.
That single sign is the whole reason verticals are chosen for low-angle work. Both
panels plot elevation angle \(\psi\) upward from the horizon at the right-hand
edge.}
\label{fig:ground-image}
\end{figure}

\subsection{The Folded Dipole's Factor of Four}
\label{sec:foldeddipole}
A folded dipole is a half-wave dipole whose ends are bent back and joined, making a
narrow loop fed at the middle of one side. Its feed impedance is famously about four
times a plain dipole's, and the reason is \cref{sec:feedpoint}'s power argument plus
\cref{sec:stubs}'s stub.

Decompose the excitation into two modes. In the \emph{transmission-line} mode the two
conductors are driven in opposition, which makes each half a two-wire line a quarter
wave long, shorted at the far end---and \cref{sec:stubs} says a shorted quarter-wave
stub has infinite input impedance. That mode therefore draws no current and radiates
nothing. Everything happens in the \emph{antenna} mode, where both conductors carry
current in the same direction and the pair radiates like one fat dipole.

Now count current. The feed injects \(I\), and in the antenna mode that current
appears in \emph{both} conductors, so the structure radiates as a dipole carrying
\(2I\). The radiated power is set by the total current and the ordinary dipole's
radiation resistance, \(P=(2I)^2R_d\), while the feed point sees \(P=I^2Z\). Equate:
\[
Z=4R_d\approx4\times\SI{73}{\ohm}=\SI{292}{\ohm}\approx\SI{300}{\ohm}.
\]
So the factor of four is a current-doubling squared, exactly as in
\cref{sec:feedpoint}---and it explains why folded dipoles are fed with
\SI{300}{\ohm} twin-lead, and why a folded dipole is the driven element of choice
when a \(4{:}1\) step-up is wanted for free.

\subsection{Terminated Antennas: Choosing \texorpdfstring{\(\Gamma_L=0\)}{Gamma L = 0}}
\label{sec:travelingwave}
A rhombic, a terminated V, or a terminated long wire carries a resistor at its far end
instead of an open circuit, and everything about such an antenna follows from
\cref{sec:reflectionsection}'s single formula. Choose the terminating resistor equal to
the structure's characteristic impedance and \(\Gamma_L=0\): no reflected wave, so no
standing wave, so the current is a \emph{traveling} wave running one way along the
wire. Three consequences, all of them the reflection coefficient's:
\begin{itemize}
\item The pattern is \textbf{unidirectional}. A standing wave has current flowing both
ways and radiates both ways; a traveling wave has a direction, and the pattern
inherits it.
\item The bandwidth is \textbf{broad}. Resonance is a standing-wave phenomenon---it is
where the reflected wave returns in phase---so an antenna with no reflection has no
resonance to be narrow about. Its behavior changes only as slowly as \(\Gamma_L=0\)
degrades.
\item The efficiency is \textbf{poor}. The power that would have been reflected is
instead dissipated in the terminator, and by \cref{sec:antennaeff} that resistor is
pure \(R_{\text{loss}}\). Half the input power is a typical figure. \emph{You are
buying pattern and bandwidth with efficiency}, which is the same currency
\cref{sec:antennaeff} was spending.
\end{itemize}


\section{Matching an Antenna to a Line}
Finally, the hardware. Every scheme below is \cref{sec:lnetwork}'s network or
\cref{sec:stubs}'s stub in antenna clothing, and the last one is a power divider
built from the same quarter-wave sections.

\subsection{The Named Antenna Matches}
\label{sec:antennamatches}
Antenna books present four matching schemes under proper names---hairpin (or beta),
gamma, delta, and stub---and they look like four separate tricks. They are not. Every
one is an L-network from \cref{sec:lnetwork} whose elements are built from things this
chapter has already derived: stub reactances (\cref{sec:stubs}) and the feed-point
tap of \cref{sec:feedpoint}. No new mathematics is needed, and one of them turns a
memorized exam fact into a consequence.

\textbf{The hairpin, and why the element must be short.} A Yagi's driven element,
pulled down by the parasitic elements coupled to it, presents a resistance well
\emph{below} \SI{50}{\ohm}---call it \SI{25}{\ohm}. Matching that up needs
\cref{sec:lnetwork}'s network with the shunt element on the \SI{50}{\ohm} side:
\[
Q=\sqrt{\frac{50}{25}-1}=1,
\qquad
|X_{\text{series}}|=QR_{\text{lo}}=\SI{25}{\ohm},
\qquad
|X_{\text{shunt}}|=\frac{R_{\text{hi}}}{Q}=\SI{50}{\ohm}.
\]
A \emph{hairpin} is a shorted stub---two parallel rods with a strap across the far
end---placed in shunt right at the feed point. \Cref{sec:stubs} gives its reactance as
\(\jj Z_0\tan\beta\ell\), and for any length under \(\lambda/4\) the tangent is
positive, so \emph{a hairpin is an inductor}. It supplies the shunt element.

Now the consequence. \Cref{sec:lnetwork} showed the two L-network reactances must be of
\emph{opposite} type, because the transformed reactance has to be cancelled. The shunt
element is inductive, so the series element must be \emph{capacitive}---and nobody
installs a capacitor. Instead the driven element is cut \emph{electrically shorter
than a half wave}, which by \cref{sec:antennalength} puts it below resonance, where
\cref{fig:antenna-model}'s series-RLC model is capacitive. \emph{``The driven element
must be capacitive'' is not a rule about hairpins; it is the L-network demanding an
opposite-type series element, with the antenna's own detuning supplying it.} The
required \SI{25}{\ohm} of capacitive reactance says exactly how much shorter.

That also fixes the hardware. With hairpin rods of
\(Z_0\approx\SI{300}{\ohm}\), \(\jj Z_0\tan\beta\ell=\jj\SI{50}{\ohm}\) needs
\(\beta\ell=\arctan(50/300)=9.5^\circ\), so \(\ell\approx0.026\lambda\)---about
\SI{5.5}{\centi\meter} at \SI{144}{\mega\hertz}. A few inches of wire, as observed.
And because the hairpin bridges the two feed terminals, both halves of the driven
element must float, which is why a hairpin-matched Yagi has its driven element
\emph{insulated from the boom} while the other schemes need not.

\textbf{Gamma and delta: tap up instead.} These take the other route to the same
place. Rather than transform the centre resistance with reactances, they move the
connection \emph{outward} along the element and use \cref{sec:feedpoint}'s
\(R(z)=R_{\text{centre}}/\cos^2\beta z\) to find a point where the resistance already
equals what the line wants. For a \SI{25}{\ohm} element, \(\lambda/8\) out doubles it
to exactly \SI{50}{\ohm}. A folded driven element (\cref{sec:foldeddipole}) reaches a
similar step-up by a different route. A \emph{delta} match does this symmetrically, fanning the
feed to two points either side of centre; a \emph{gamma} match does it
\emph{unbalanced}, bonding the coax shield to the centre and tapping the centre
conductor out to one side---which is precisely why the gamma is the one used to
shunt-feed a grounded tower at its base, since one side of the feed is already at
ground.

But a tap is not only a resistance change: the conductor running out to the tap and
back is a loop, so it adds series \emph{inductance}. Something must cancel it, and
here the antenna cannot help---the tap point is chosen for resistance, not reactance.
Hence the series capacitor that every gamma match carries. Its purpose is not DC
isolation and not impedance transformation; it \emph{cancels the unwanted inductive
reactance} the tap arrangement introduced. Same L-network logic as the hairpin, with
the roles of the two elements exchanged.

\textbf{Stub match: the general case.} A shorted or open stub placed in parallel at or
near the feed point, its length and position both adjustable, is simply
\cref{sec:stubs} used as a two-degree-of-freedom matching network (and
\cref{sec:smithuse} is how the chart finds those two numbers by eye)---the position
choosing what impedance the line presents at that plane, the stub length supplying the
susceptance that cancels its imaginary part. It needs no antenna modification at all,
which is why it works on anything and why the others exist only as tidier packages for
specific geometries.

\begin{exambox}
Group E9E names six of these, and every answer is above (its seventh, the power
divider, is \cref{sec:wilkinson}). The \emph{beta or hairpin}
match is the one needing an insulated driven element, and it needs that element
\emph{capacitive---electrically shorter than a half wave}. The \emph{gamma} match is
the one connecting shield to centre and conductor a fraction of a wavelength to one
side, the one that shunt-feeds a grounded tower, and the one whose series capacitor
exists \emph{to cancel unwanted inductive reactance}. The \emph{stub} match is the
short length of line in parallel at or near the feed point. And E9E06's Q-section is
\cref{sec:quarterwave} arithmetic: \(\sqrt{50\times100}=\SI{70.7}{\ohm}\), so of the
choices offered the \SI{75}{\ohm} cable is the one to use.
\end{exambox}

\subsection{The Wilkinson Divider: Two Modes, Three Requirements}
\label{sec:wilkinson}
Splitting power between two antennas sounds easy---join the cables---but a bare
junction fails all three things a divider must do at once: split equally, present
\(Z_0\) at the input, and keep the two outputs from talking to each other. The
Wilkinson divider achieves all three with two quarter-wave sections and one resistor,
and the reason is a symmetry argument of exactly the kind \cref{sec:pushpull} used.

The circuit: from the input, two \(\lambda/4\) sections of characteristic impedance
\(Z_0\sqrt2\) run to the two output ports, and a resistor of \(2Z_0\) bridges the
outputs. For \(Z_0=\SI{50}{\ohm}\) that is \SI{70.7}{\ohm} sections and a
\SI{100}{\ohm} resistor. Now split any excitation of the two outputs into an
\emph{even} part (equal and in phase) and an \emph{odd} part (equal and opposite), and
analyze each.

\emph{Even mode.} Both outputs are at the same potential, so \emph{no current flows in
the bridging resistor}---it might as well not be there. Each \(\lambda/4\) section is
then a transformer (\cref{sec:quarterwave}) loaded by \(Z_0\), presenting
\((Z_0\sqrt2)^2/Z_0=2Z_0\) at the input. Two of those in parallel give \(Z_0\). The
input is matched, and the power divides equally by symmetry.

\emph{Odd mode.} The two outputs are equal and opposite, so the midpoint of the
bridging resistor is a \emph{virtual ground}, and so is the input node. Each output
therefore sees half the resistor, \(Z_0\), in parallel with a \(\lambda/4\) section
that is \emph{shorted} at the input plane---and a shorted quarter wave looks like an
open circuit (\cref{sec:stubs}). The open contributes nothing, so the output port sees
exactly \(Z_0\). That is the isolation: any odd-mode signal injected at one output is
absorbed in the resistor and never reaches the other output or the input.

So the resistor is invisible to the signal you want and is the only thing the signal
you do not want can see. \emph{The two modes do not interact, which is why one
component can satisfy requirements that look contradictory.} Note also what the
resistor is not doing: in normal operation it dissipates nothing, so it need only be
rated for the imbalance, not for the through power---until one antenna is damaged, at
which point it absorbs the difference and tells you so by getting hot.

\begin{exambox}
E9E asks what a Wilkinson divider is for, and the answer is the three requirements
above stated together: \emph{to divide power equally between two \SI{50}{\ohm} loads
while maintaining \SI{50}{\ohm} input impedance}. The isolation is the part the answer
does not mention and the part the resistor is actually there for.
\end{exambox}


\begin{figure}[htbp]
\centering
\begin{circuitikz}[scale=1]
\draw (0,0) node[left]{\(\circ\)} to[R,l=\(R_{\text{rad}}\)] (2.2,0)
      to[L,l=\(L\)] (4.4,0)
      to[C,l=\(C\)] (6.4,0) node[right]{\(\circ\)};
\end{circuitikz}
\caption{Idealized antenna as a series-RLC one-port: radiation resistance
\(R_{\text{rad}}\) in series with the reactance that resonates at
\(\omega_0=1/\sqrt{LC}\). The two terminals are the feed point a NanoVNA---a
pocket-sized vector network analyzer (VNA)---connects to.}
\label{fig:antenna-model}
\end{figure}

\section{Worked Example}
A \(\SI{100}{\ohm}\) load terminates a \(\SI{50}{\ohm}\) line:
\[
\Gamma=\frac{100-50}{100+50}=\frac{1}{3},
\qquad
\mathrm{SWR}=\frac{1+1/3}{1-1/3}=2,
\qquad
RL=-20\log_{10}(1/3)\approx \SI{9.54}{\decibel}.
\]
A quarter-wave section would need
\(Z_t=\sqrt{(50)(100)}=\SI{70.7}{\ohm}\) to remove the mismatch, giving
\(Z_{\mathrm{in}}=Z_t^2/Z_L=70.7^2/100=\SI{50}{\ohm}\).

\Cref{fig:line-worked} shows both halves of this example, and each panel makes a point
the algebra alone leaves abstract.

The left panel is the standing wave on the \emph{unmatched} line. Since
\(\Gamma=+1/3\) is real and positive, the reflected wave arrives in phase with the
incident one right at the load, so \(x=0\) is a voltage \emph{maximum} of
\(1+|\Gamma|=1.33\); a quarter wavelength back the round trip has added
\(180^\circ\) and the waves subtract, giving the minimum \(1-|\Gamma|=0.67\). Their
ratio is the SWR of \(2\) computed above---so ``SWR${}=2$'' means literally that a
voltmeter slid along this line would read twice as much in some places as in others.
Note the spacing: maxima repeat every \emph{half} wavelength, and a maximum is a
\emph{quarter} wavelength from the neighboring minimum.

The right panel sweeps frequency with the quarter-wave section in place. The match is
exact only at \(f_0\), where the section is precisely \(\lambda/4\); away from it
\(\beta\ell\) is no longer \(\pi/2\), the impedance inversion is imperfect, and SWR
climbs. This is what ``narrowband'' means quantitatively, and it is the reason a
quarter-wave transformer is a single-band device while the lumped L-network of
\cref{ch:filters}---which can be retuned---is not.

\begin{figure}[htbp]
\centering
\includegraphics[width=\textwidth]{line_worked}
\caption{The worked example. \emph{Left:} voltage standing-wave envelope
\(|V(x)|=|V^{+}||1+\Gamma\ee^{-\jj2\beta x}|\) on the unmatched line, whose
maximum-to-minimum ratio is the SWR of \(2\); maxima recur every \(\lambda/2\) and
sit \(\lambda/4\) from the minima. \emph{Right:} SWR versus frequency once the
\(\SI{70.7}{\ohm}\) quarter-wave section is inserted (designed for
\(f_0=\SI{14.1}{\mega\hertz}\)). The match is perfect only at \(f_0\) and degrades
either side of it, since the section's electrical length is proportional to
frequency.}
\label{fig:line-worked}
\end{figure}
